\begin{lemma}[A strict augmentation route starts with a colored boundary]
Let \(A\) and \(B\) be tree-sequences for the same finite loopless
multigraph and distinguished first edge, and let \(\varphi\) be a fixed
partial edge coloring.  Suppose that \(B\) is obtained from \(A\) by a
\(\varphi\)-\noderef{TashkinovAugmentationSequence} and that
\(|A|<|B|\).  Then there is a color
\(c\in\noderef{TreeMissingColors}(\varphi,A)\) for which the exact colored
boundary set
\(\noderef{BoundaryColorEdges}(G,\varphi,V(A),c)\) is nonempty.
\end{lemma}
\begin{proof}
By \noderef{TashkinovAugmentationSequence}, the given route is a finite
reflexive-transitive chain of one-step
\(\varphi\)-\noderef{TashkinovAugmentation}s.  The strict length inequality
rules out the empty chain.  Isolate its first step, say \(A\to C\).

By \noderef{TashkinovAugmentation}, \(C\) has length \(|A|+1\), \(A\) is
its exact ordered initial segment, and its new last edge \(f\) has a color
\(c\) belonging to \noderef{TreeMissingColors}\((\varphi,A)\).  It remains
only to verify that this exact edge identity \(f\) is a boundary edge of
\(V(A)\).

Write \(j=|A|\) for the index of \(f\) in \(C\).  The attachment clause in
\noderef{TreeSequence} supplies an index \(a\leq j\) such that the ends of
\(f\) are the new vertex \(C.y_{j+1}\) and the earlier vertex \(C.y_a\).
Because \(A\) is an exact ordered
\noderef{TreeSequenceInitialSegment} of \(C\), one has
\(C.y_a=A.y_a\), so this earlier endpoint belongs to
\(\noderef{TreeSequenceVertexSet}(A)\).

The new endpoint \(C.y_{j+1}\) does not belong to \(V(A)\).  Otherwise it
would equal \(A.y_b=C.y_b\) for some \(b\leq |A|=j\), contradicting the
injectivity of the vertex map of \(C\), since \(b<j+1\).  Thus \(f\) has
one endpoint in \(V(A)\) and its other endpoint outside \(V(A)\), so
\(f\in\noderef{BoundaryEdges}(G,V(A))\).  The first-step augmentation also
gives \(\varphi(f)=c\).  By the definition of
\noderef{BoundaryColorEdges}, it follows that
\(f\in\partial_{\varphi,c}(V(A))\).  Hence that exact colored boundary
set is nonempty, witnessed by \(f\), and the chosen \(c\) has all the
required properties.
\end{proof}
